\section{Introduction}
\label{sec:introduction}

Open-weight mixture-of-experts (MoE) models decouple total model capacity from
the computation performed for each token. A router selects a small subset of
experts, so a model with tens or hundreds of billions of parameters may execute
only a few billion active parameters at a time. That sparsity makes local
serving possible, but it does not make the inactive weights disappear. A local
runtime must still place the complete expert pool and move requested experts
into execution memory. It must also decide which cached state should survive
when the machine becomes busy.

FreeToken addresses this problem on heterogeneous CUDA systems by combining a
host-resident expert pool, a GPU expert cache, double-buffered prefill, and a
bandwidth-calibrated choice between CPU execution and PCIe transfer
\cite{yang2026freetoken}. KTransformers similarly demonstrates that CPU-GPU
hybrid execution can make large MoE checkpoints practical on commodity
machines \cite{chen2025ktransformers}. These systems exploit a resource model
with distinct host and device memories. Apple Silicon changes that model. The
CPU and GPU share physical memory, and Metal shared buffers do not cross a PCIe
boundary. Expert slots, immutable file-backed weights, prompt state, the
filesystem cache, the display stack, desktop applications, compression, and
swap all compete in the same pool.

This shared pool creates a failure mode that is invisible to cache-local
metrics. We first characterize whether cache hit rate tracks end-to-end service
performance. Increasing explicit expert capacity can raise the hit rate yet
evict useful file-backed pages, increase compression, and reduce decode speed.
Our preregistered fixed-cache sweep demonstrates that tension. Increasing Qwen
from 16 to 96 expert slots per layer raises median hit rate from 48.98\% to
90.13\%, while median decode throughput falls from 11.330 to 10.120\tps. Peak
compressor growth also increases by roughly 5\gib. The paired slowdown is
10.54\%, below the frozen 15\% threshold required to authorize our original
cache-size controller. We retain that negative result and abandon the
controller rather than lower its gate.

The fixed-cache curve suggests a different mechanism. The runtime benefits from
the logical identity and reuse information of a large cache, but it need not
claim permanent physical residency for every cold slot. Apple exposes
purgeable Metal resources for exactly this cache-like relationship: an
application may mark a resource volatile, the operating system may reclaim its
contents, and the application must relock and recreate or reload the resource
before access \cite{apple2026purgeable,apple2026metalfootprint}. Applying that
contract to model weights is not automatic. A cache mapping can outlive the
physical contents it names; a GPU command may still reference a resource when a
state transition is attempted; and externally allocated memory wrapped by
Metal may remain physically owned even when the wrapper is marked purgeable.

We present \system to resolve those correctness and integration problems. Each
expert slot is allocated through Metal and retains a stable CPU pointer for
positional reads. The cache tracks each slot as nonvolatile, volatile, or
empty. At safe boundaries it
keeps the hottest slots nonvolatile and offers the cold remainder to macOS. A
future logical hit is provisional until a nonvolatile relock confirms that its
contents survived. An empty return clears identity and recency state and enters
the existing exact miss path. The GPU sees either validated cached bytes or a
fully reloaded expert; discarded contents cannot reach a command encoder.

This paper makes four contributions:

\begin{itemize}
  \item It characterizes the mismatch between expert-cache hit rate and
  end-to-end service performance in one Apple unified-memory system, including
  a preregistered negative result that stops an earlier controller claim.
  \item It presents a reload-correct state machine for Metal-owned expert slots,
  with explicit validity, publication, and GPU-synchronization invariants.
  \item It integrates selective purgeability at the actual prefill-layer and
  decode-token boundaries, avoiding the full physical prefill peak that defeats
  a post-prefill-only policy.
  \item It evaluates the committed implementation against untouched small- and
  large-cache upstream binaries on Qwen3.6 and Gemma 4, with balanced held-out
  repetitions, physical-footprint telemetry, exact-recovery receipts, negative
  alternatives, immutable artifacts, and explicit claim limits.
\end{itemize}

The strongest supported conclusion is intentionally scoped. On the measured M1
Max workloads, \system improves the throughput-footprint frontier: it dominates
fixed-96 with higher throughput and lower footprint, while providing more
throughput and locality than fixed-16 at a higher footprint. One host and two
architectures do not establish an all-Mac result. Cross-chip reproduction,
energy measurement, longer contexts, concurrency, and independent operation
remain required.
